\documentclass[a4paper,10pt]{article}

\usepackage{latexsym}
\usepackage[T1]{fontenc}
\usepackage[utf8]{inputenc}
\usepackage[empty]{fullpage}
\usepackage{titlesec}
\usepackage{marvosym}
\usepackage[usenames,dvipsnames]{color}
\usepackage{verbatim}
\usepackage{enumitem}
\usepackage[hidelinks]{hyperref}
\usepackage{fancyhdr}
\usepackage[french]{babel}
\usepackage{tabularx}
\input{glyphtounicode}

\pagestyle{fancy}
\fancyhf{}
\fancyfoot{}
\renewcommand{\headrulewidth}{0pt}
\renewcommand{\footrulewidth}{0pt}

\addtolength{\oddsidemargin}{-0.5in}
\addtolength{\evensidemargin}{-0.5in}
\addtolength{\textwidth}{1in}
\addtolength{\topmargin}{-.7in}
\addtolength{\textheight}{1.4in}

\urlstyle{same}
\raggedbottom
\raggedright
\setlength{\tabcolsep}{0in}

\titleformat{\section}{
  \vspace{-8pt}\scshape\raggedright\large
}{}{0em}{}[\color{black}\titlerule \vspace{-5pt}]

\pdfgentounicode=1

% --- CUSTOM COMMANDS ---
\newcommand{\resumeItem}[1]{
  \item\small{#1 \vspace{0pt}}
}

\newcommand{\resumeSubheading}[4]{
  \vspace{0pt}\item
    \begin{tabular*}{0.97\textwidth}[t]{l@{\extracolsep{\fill}}r}
      \textbf{#1} & #2 \\
      \textit{\small#3} & \textit{\small #4} \\
    \end{tabular*}\vspace{-5pt}
}

\newcommand{\resumeProjectHeading}[2]{
    \vspace{1pt}\item
    \begin{tabular*}{0.97\textwidth}{l@{\extracolsep{\fill}}r}
      \small#1 & #2 \\
    \end{tabular*}\vspace{-3pt}
}

\newcommand{\resumeSubItem}[1]{\resumeItem{#1}\vspace{-2pt}}
\renewcommand\labelitemii{$\vcenter{\hbox{\tiny$\bullet$}}$}
\newcommand{\resumeSubHeadingListStart}{\begin{itemize}[leftmargin=0.15in, label={}]}
\newcommand{\resumeSubHeadingListEnd}{\end{itemize}}
\newcommand{\resumeItemListStart}{\begin{itemize}[label=\textbullet]}
\newcommand{\resumeItemListEnd}{\end{itemize}\vspace{-4pt}}

\begin{document}

%----------EN-TÊTE----------
\begin{center}
    \textbf{\Huge \scshape Théo Phan} \\ \vspace{2pt}
    \large Stagiaire Ingénieur Logiciel \quad $|$ \quad Disponible été 2027 \\ \vspace{2pt}
    \small
    \href{mailto:theo.phan.quoc.huy@gmail.com}{\underline{theo.phan.quoc.huy@gmail.com}} $|$
    \href{tel:+33613620366}{\underline{+33 6 13 62 03 66}} $|$
    \href{https://linkedin.com/in/theophanquochuy/}{\underline{LinkedIn}} $|$
    \href{https://github.com/Aer-3888}{\underline{GitHub}} $|$
    \href{https://portfolio-theo.pages.dev/}{\underline{Portfolio}} $|$
    ouvert à la mobilité
\end{center}

%-----------EXPÉRIENCE-----------
\section{Expérience}
  \resumeSubHeadingListStart

    \resumeSubheading
      {Stagiaire Ingénieur Logiciel, Machine Learning appliqué}{Avril -- Juin 2025}
      {FPT Telecom -- détection de défauts en temps réel sur câblage fibre optique}{Hanoï, Vietnam}
      \resumeItemListStart
        \resumeItem{Conception et livraison d'un \textbf{pipeline Python de bout en bout} (collecte, annotation, augmentation, entraînement, inférence avec \textbf{PyTorch / YOLOv8}) transformant une inspection visuelle manuelle en un workflow de \textbf{détection en temps réel} automatisé, déployé dans un outil mobile interne utilisé sur site par les techniciens}
        \resumeItem{Comparaison de deux architectures de modèles sur des métriques objectives (mAP, précision/rappel, latence) pour guider la décision de production, et mise en place de \textbf{contrôles qualité des données} détectant lacunes et déséquilibre de classes avant chaque entraînement}
        \resumeItem{Rédaction d'une \textbf{documentation technique} complète pour une passation propre aux équipes d'ingénierie internes}
      \resumeItemListEnd

    \resumeSubheading
      {Founding Engineer \& Tech Lead, Waiki}{2025 -- Présent}
      {Startup en amorçage -- matériel NFC en chambre et application compagnon multiplateforme}{Hybride}
      \resumeItemListStart
          \resumeItem{Conception et \textbf{mise en production} (iOS/Android) de l'application compagnon, aujourd'hui \textbf{1{,}425+ unités vendues}, avec le backend et la couche de données (Firebase, SQLite, \textbf{MVVM}) pour un état appareil/utilisateur fiable et synchronisé}
          \resumeItem{Pilotage de la livraison, de besoins ambigus aux fonctionnalités livrées, en dirigeant une équipe de 3 en \textbf{Agile/Scrum} (user stories, Kanban, sprint planning) avec des résultats mesurables à chaque sprint}
      \resumeItemListEnd

    \resumeSubheading
      {Développeur Backend (Bénévole), Gnut06}{2026 -- Présent}
      {Association utilisant la VR/AR pour l'inclusion des personnes en situation de handicap}{À distance}
      \resumeItemListStart
        \resumeItem{Développement de nouvelles fonctionnalités et de la couche de données \textbf{Doctrine / MySQL} d'une plateforme web \textbf{Symfony 7.1 / PHP} (\textbf{Docker}, Webpack Encore), dans un workflow Git avec revues de code}
        \resumeItem{\textbf{Refactorisation de code legacy} vers une base plus propre, plus sûre et plus maintenable, réduisant la dette technique}
      \resumeItemListEnd

  \resumeSubHeadingListEnd

%-----------PROJETS-----------
\section{Projets \& Open Source}
  \resumeSubHeadingListStart

    \resumeProjectHeading
      {\textbf{\href{https://github.com/Aer-3888/Notes_insa}{\underline{Notes INSA, Analytics Étudiant Distribué}}}}{\emph{TypeScript, Cloudflare Workers, D1 (SQL), Flutter}}
      \resumeItemListStart
        \resumeItem{Développement et publication d'une app mobile \textbf{Flutter} permettant aux étudiants de \textbf{se situer} par rapport aux statistiques de la promotion en direct (percentile, moyenne, médiane), calculées côté client pour la confidentialité et utilisées par de vrais étudiants}
        \resumeItem{Conception du \textbf{backend serverless edge} (\textbf{TypeScript} sur \textbf{Cloudflare Workers}, SQL sur \textbf{D1}) ingérant, dédupliquant et agrégeant en temps réel les soumissions de toute une promotion, un système \textbf{distribué et scalable} de stockage et de requêtes}
      \resumeItemListEnd

    \resumeProjectHeading
      {\textbf{\href{https://github.com/Aer-3888/iaj_evades}{\underline{IAJ Evades, Agent d'Apprentissage par Renforcement}}}}{\emph{Rust, TypeScript, DQN, WebSockets}}
      \resumeItemListStart
        \resumeItem{Développement d'un moteur de jeu en \textbf{Rust} et d'un framework d'entraînement \textbf{Deep Q-Network} avec \textbf{entraînement par batch multi-threadé} et focalisation sur les mauvais seeds pour la robustesse de l'agent}
        \resumeItem{Réalisation d'un dashboard \textbf{TypeScript} diffusant les \textbf{métriques d'entraînement en temps réel via WebSockets}, avec réglage des paramètres à la volée et checkpointing JSON pour des runs reprenables et tolérants aux pannes}
      \resumeItemListEnd

    \resumeProjectHeading
      {\textbf{\href{https://github.com/Aer-3888/GardeFou}{\underline{GardeFou, Laboratoire de Sécurité LLM}}}}{\emph{Python, RAG, Prompt Injection, Ollama, Chroma}}
      \resumeItemListStart
        \resumeItem{Construction d'un assistant \textbf{RAG} volontairement vulnérable sur des documents financiers, puis cycle \textbf{red team / blue team} pour démontrer les attaques par \textbf{prompt injection} et les garde-fous qui les stoppent}
        \resumeItem{Vérification que chaque attaque échoue après durcissement et mapping de chaque découverte à un \textbf{référentiel de sécurité LLM} reconnu, avec preuves avant/après}
      \resumeItemListEnd

  \resumeSubHeadingListEnd

%-----------COMPÉTENCES-----------
\section{Compétences}
\begin{itemize}[nosep,leftmargin=0.15in,label=\textbullet]
    \resumeItem{\textbf{Langages de programmation :} Java, Python, C++, C\#, TypeScript/JavaScript, Rust, SQL.}
    \resumeItem{\textbf{Fondamentaux CS :} conception orientée objet, structures de données \& algorithmes, analyse de complexité, concurrence \& multithreading, systèmes d'exploitation, systèmes distribués.}
    \resumeItem{\textbf{IA / ML :} PyTorch, TensorFlow, réseaux de neurones (CNN, DQN), vision par ordinateur (YOLOv8, OpenCV), Scikit-learn, ONNX Runtime.}
    \resumeItem{\textbf{Backend \& Data :} API REST, WebSockets, serverless (Cloudflare Workers), Node.js, bases relationnelles (SQL, MySQL, D1), MongoDB.}
    \resumeItem{\textbf{Outils :} Docker, CI/CD (GitHub Actions), Git/GitLab, Linux.}
    \resumeItem{\textbf{Langues :} Français (natif), Vietnamien (natif), Anglais (professionnel, C1).}
\end{itemize}

%-----------FORMATION-----------
\section{Formation}
  \resumeSubHeadingListStart
    \resumeSubheading
      {INSA Rennes, Diplôme d'Ingénieur (niveau Master) en Informatique}{2025 -- 2028}
      {Spécialisation IA. Cours : algorithmes \& structures de données, systèmes d'exploitation, probabilités}{Rennes, France}
    \resumeSubheading
      {IUT de Nantes, BUT Informatique}{2023 -- 2025}
      {Parcours : Réalisation d'Applications, Conception et Validation}{Nantes, France}
  \resumeSubHeadingListEnd

%-----------ACTIVITÉS-----------
\section{Activités}
\begin{itemize}[nosep,leftmargin=0.15in,label=\textbullet]
    \resumeItem{\textbf{CTF} (Équipe Crabe) : 2\textsuperscript{e} place IUT CTF, 71\textsuperscript{e}/757 au classement mondial du \href{https://ctftime.org/event/3105}{\underline{Midnight Flag CTF 2026}}.}
    \resumeItem{\textbf{AI for Business Hackathon 2026} : projet en équipe appliquant des solutions IA à des problématiques business.}
    \resumeItem{\textbf{Photographie \& guitare} : pratique personnelle autodidacte depuis 2 ans.}
\end{itemize}

\end{document}
